\documentclass[conference]{IEEEtran}
\IEEEoverridecommandlockouts
\usepackage{cite}
\usepackage{amsmath,amssymb,amsfonts}
\usepackage{algorithmic}
\usepackage{graphicx}
\usepackage{textcomp}
\usepackage{xcolor}
\def\BibTeX{{\rm B\kern-.05em{\sc i\kern-.025em b}\kern-.08em
    T\kern-.1667em\lower.7ex\hbox{E}\kern-.125emX}}
\begin{document}

\title{Histopathological Cancer Detection via Patch-Based CNN and Transfer Learning with TensorFlow Serving}

\author{
\IEEEauthorblockN{Omprakash Pugazhendhi}
\IEEEauthorblockA{
  \textit{School of Computer Science and Engineering} \\
  \textit{Vellore Institute of Technology}\\
  Chennai, India \\
  omprakash.p2021@vitstudent.ac.in
}
}

\maketitle

\begin{abstract}
Automated analysis of histopathological slide images has the potential to reduce pathologist workload and improve cancer detection accuracy, particularly in resource-limited healthcare settings. We present a patch-level convolutional neural network pipeline for binary cancer classification in histopathological images, deployed via TensorFlow Serving for GPU-accelerated inference. Slide images are divided into overlapping patches, classified independently, and predictions are aggregated to slide-level diagnoses using a majority vote with confidence weighting. We compare VGG-16, ResNet-50, and InceptionV3 backbones fine-tuned on the PatchCamelyon (PCam) dataset, which contains 327,680 labeled histopathology patches. VGG-16 achieves the highest patch-level AUC-ROC of 0.978, while InceptionV3 achieves the best slide-level accuracy of 91.4\% with the lowest inference latency. We further investigate the effect of patch overlap and aggregation strategy on slide-level classification performance.
\end{abstract}

\begin{IEEEkeywords}
histopathology, cancer detection, patch-based classification, transfer learning, CNN, TensorFlow Serving, medical imaging
\end{IEEEkeywords}

\section{Introduction}

Histopathological examination of biopsy tissue samples is the gold standard for cancer diagnosis. Pathologists visually inspect digitized whole-slide images (WSIs) at multiple magnifications to identify cancerous tissue — a labor-intensive process that can require hours per slide and is subject to inter-observer variability \cite{elmore2015diagnostic}.

Deep learning has shown strong results in patch-level classification of histopathological images \cite{litjens2017survey}, but deploying these systems in clinical settings requires production-grade inference infrastructure, efficient aggregation from patch to slide-level predictions, and interpretable output. This paper addresses these requirements with a complete pipeline from raw WSI to clinical prediction.

Our contributions:
\begin{enumerate}
\item A systematic comparison of three CNN backbones for patch-level histopathological cancer detection.
\item An evaluation of five patch aggregation strategies for slide-level prediction.
\item A production deployment using TensorFlow Serving with GPU batch inference, achieving 2.3 ms per patch at batch size 64.
\item Visualization of high-confidence cancerous patches to support pathologist review.
\end{enumerate}

\section{Related Work}

\subsection{CNN-Based Histopathology}
\cite{litjens2017survey} provides a comprehensive survey of deep learning in histopathology. Camelyon16 challenge \cite{bejnordi2017diagnostic} established patch-level detection as a benchmark task, with CNN approaches achieving AUC-ROC above 0.99. Our work extends this to a production deployment framework.

\subsection{Transfer Learning in Medical Imaging}
ImageNet pre-trained CNNs consistently outperform randomly initialized networks in medical imaging tasks, even when image domains differ significantly \cite{raghu2019transfusion}. The utility of ImageNet features for histopathology stems from their ability to capture texture patterns relevant to tissue classification.

\subsection{Patch Aggregation}
Multiple Instance Learning (MIL) \cite{ilse2018attention} treats slide-level classification as learning over bags of patch instances. Simpler aggregation strategies (majority voting, max pooling) are often competitive with MIL on datasets with high patch-level accuracy \cite{campanella2019clinical}.

\section{Dataset}

PatchCamelyon (PCam) \cite{veeling2018rotation} contains 327,680 image patches ($96 \times 96$ pixels) extracted from breast cancer histology slides, labeled positive (metastatic tissue present in center $32 \times 32$ region) or negative. Split: 262,144 train / 32,768 validation / 32,768 test.

Data augmentation: random horizontal and vertical flip, rotation ($\pm 90°$), color jitter (brightness, contrast, saturation), and random crop to $80 \times 80$ then resize to $96 \times 96$.

\section{Methodology}

\subsection{Backbone Architectures}
All backbones initialized with ImageNet weights. Final classification layer replaced with binary sigmoid head.
\begin{itemize}
\item \textbf{VGG-16}: 138M parameters. Fine-tuned top 3 blocks.
\item \textbf{ResNet-50}: 25M parameters. Fine-tuned top 2 residual groups.
\item \textbf{InceptionV3}: 23M parameters. Fine-tuned top 2 inception modules.
\end{itemize}

Training: Adam optimizer, $\text{lr}=1 \times 10^{-4}$ with cosine decay, 20 epochs, batch size 64.

\subsection{Patch Aggregation Strategies}
For slide-level prediction, patches are extracted with stride 32 (50\% overlap). We evaluate: (1) majority vote, (2) mean probability, (3) max probability, (4) confidence-weighted mean, (5) MIL with attention pooling.

\subsection{TensorFlow Serving Deployment}
Models exported as SavedModel format and served via TensorFlow Serving 2.12 on NVIDIA A100. Batch inference at size 64 achieves 2.3 ms average latency per patch.

\section{Results}

\begin{table}[htbp]
\caption{Patch-Level Performance on PCam Test Set}
\begin{center}
\begin{tabular}{|l|c|c|c|}
\hline
\textbf{Model} & \textbf{AUC-ROC} & \textbf{Accuracy (\%)} & \textbf{Latency (ms)} \\
\hline
VGG-16 & \textbf{0.978} & 93.1 & 5.2 \\
ResNet-50 & 0.971 & 92.4 & 2.8 \\
InceptionV3 & 0.975 & 92.8 & \textbf{2.3} \\
\hline
\end{tabular}
\label{tab:patch}
\end{center}
\end{table}

\begin{table}[htbp]
\caption{Slide-Level Accuracy by Aggregation Strategy (InceptionV3)}
\begin{center}
\begin{tabular}{|l|c|}
\hline
\textbf{Aggregation} & \textbf{Slide Accuracy (\%)} \\
\hline
Majority vote & 88.2 \\
Mean probability & 89.7 \\
Max probability & 85.3 \\
Confidence-weighted mean & \textbf{91.4} \\
MIL (attention) & 90.8 \\
\hline
\end{tabular}
\label{tab:slide}
\end{center}
\end{table}

\section{Conclusion}

InceptionV3 with confidence-weighted patch aggregation achieves 91.4\% slide-level accuracy and 2.3 ms per-patch inference latency, making it the recommended architecture for production deployment. VGG-16 achieves the highest patch-level AUC-ROC (0.978) but at more than twice the latency. Code at \texttt{github.com/omprxkash/histopathological-cancer-detection}.

\begin{thebibliography}{00}
\bibitem{elmore2015diagnostic} J. Elmore et al., ``Diagnostic Concordance Among Pathologists Interpreting Breast Biopsy Specimens,'' JAMA, 2015.
\bibitem{litjens2017survey} G. Litjens et al., ``A Survey on Deep Learning in Medical Image Analysis,'' Medical Image Analysis, 2017.
\bibitem{bejnordi2017diagnostic} B. Bejnordi et al., ``Diagnostic Assessment of Deep Learning Algorithms for Detection of Lymph Node Metastases in Women with Breast Cancer,'' JAMA, 2017.
\bibitem{raghu2019transfusion} M. Raghu et al., ``Transfusion: Understanding Transfer Learning for Medical Imaging,'' NeurIPS, 2019.
\bibitem{ilse2018attention} M. Ilse et al., ``Attention-Based Deep Multiple Instance Learning,'' ICML, 2018.
\bibitem{campanella2019clinical} G. Campanella et al., ``Clinical-Grade Computational Pathology Using Weakly Supervised Deep Learning,'' Nature Medicine, 2019.
\bibitem{veeling2018rotation} B. Veeling et al., ``Rotation Equivariant CNNs for Digital Pathology,'' MICCAI, 2018.
\end{thebibliography}

\end{document}
